\documentclass[11pt,a4paper]{article}
\usepackage[T1]{fontenc}
\usepackage{lmodern,amsmath,amssymb,booktabs,microtype}
\usepackage[margin=25mm]{geometry}
\usepackage[colorlinks=true,urlcolor=blue,linkcolor=blue]{hyperref}
\newcommand{\E}{\mathbb E}
\title{Four-time Gaussian smearing of the\\dS$_3$ eikonal observer correlator}
\author{Alan Day \and Chengheng\thanks{AI research assistant (OpenAI).}}
\date{Working research note --- 31 October 2026}
\begin{document}\maketitle
\begin{abstract}
We apply four independent real Gaussian time smearings to the existing
two-conformal-probe eikonal formula. An exact change of Gaussian variables
reduces the calculation to three time variables and an impact angle.
The disconnected normalization changes along with the numerator.
We derive the fixed-order smearing factors and show that equal real
smearing does not remove the imaginary difference between the two
specified lateral contours. Numerical quadratures check the reduction
against the original four insertion coordinates. All statements concern
the specified eikonal approximation, not a completed quantum-gravity
state or a finite-$G$ holographic dual.
\end{abstract}
\section{The comparison target}
Set the de Sitter radius to one, $m>0$, $\gamma=4Gm\in[0,1)$,
and $\alpha=2\pi(1-\gamma)$. Retain the CSTY fixed-average-mass
mixed representation and its two contour prescriptions \cite{csty}.
There are two distinguishable conformally coupled probes,
$\nu=\nu'=i/2$, with mass squared $3/4$. This is a complementary-series
special case. The preceding all-particle construction assumed real
principal-series $p>0$; we do not silently analytically continue that
Hilbert-space theorem to the present model.

The four independent real displacements $d_i$ are centered Gaussians
of variance $\sigma^2$, with fixed finite $\sigma>0$:
\[
 (\theta_1,\theta_2,\theta_3,\theta_4)
 =(-\pi/2+i(t+d_1),\,id_2,\,\pi/2+i(t+d_3),\,\pi+id_4).
\]
Each operator is smeared with the normalized kernel
$(2\pi\sigma^2)^{-1/2}\exp[-d^2/(2\sigma^2)]$.
Imaginary separations stay fixed. This is not the mass-threshold
tracial convolution in CSTY Appendix F, nor a replacement of all
four shifts by a single shift of $t$.

CLP's proposed finite-parameter insertions retain an undetermined
normalization beyond the leading probe limit \cite{clp}.
Consequently, the present calculable target is the semiclassical
eikonal correlator, not the bare tetrahedron amplitude or its
unfixed physical completion.

\newpage
\section{Reduction from the momentum integral}
For $0\le\phi\le\pi$ the existing shock profile is
\[
 h(\phi)=\frac{H_\gamma(\phi)}m,\qquad
 H_\gamma(\phi)=-\frac{\cos[(1-\gamma)(\pi-\phi)]}
                         {2\,\operatorname{sinc}(\pi\gamma)} ,
 \quad \operatorname{sinc}x=\frac{\sin x}x .
\]
Angle brackets below mean $\pi^{-1}\int_0^\pi d\phi$.
Define
\[
 X=\frac{d_1-d_3}{2},\quad Y=\frac{d_2-d_4}{2},\quad
 Z=\frac{d_1+d_3-d_2-d_4}{2},\quad T=t+Z,\quad C=\cosh X\cosh Y.
\]
Joint Gaussian covariance gives independent
$X,Y\sim N(0,\sigma^2/2)$ and $Z\sim N(0,\sigma^2)$.
The unused common center integrates to one.

In CSTY (G.2), let $A=e^{-i\theta_4}-e^{-i\theta_2}$,
$B=e^{i\theta_1}-e^{i\theta_3}$ and
$P=\exp[i(\theta_1+\theta_3-\theta_2-\theta_4)/2]$.
Writing $s_W=(d_1+d_3)/2$ and $s_V=(d_2+d_4)/2$ gives
\[
 A=-2e^{s_V}\cosh Y,\quad
 B=-2ie^{-t-s_W}\cosh X,\quad P=-ie^{-T}.
\]
The sector integral after $p_-=-v$ is
\[
 \frac{P}{4\alpha^2}
 \left\langle\int_0^\infty\frac{e^{-iBv}}{i(A-hv)}\,dv\right\rangle .
\]
Rescaling $v$ by $iB=2e^{-t-s_W}\cosh X$ gives a disconnected factor
$F_{\rm disc}(X,Y)=1/(16\alpha^2C)$ and the argument $he^T/(4C)$.
With the sector sum and lateral prescriptions of CSTY
(2.25)--(2.26), (2.33)--(2.35), define
\begin{align}
 R_\eta(T,C)&=\left\langle
 J_\eta\left(\frac{he^T}{4C}\right)
 +J_{-\eta}\left(\frac{he^{-T}}{4C}\right)-1\right\rangle,\label{R}\\
 J(c)&=\int_0^\infty\frac{e^{-x}}{1+cx}\,dx
      =\frac{e^{1/c}E_1(1/c)}c,\quad c>0,\nonumber\\
 J_\eta(-1/a)&=ae^{-a}\operatorname{Ei}(a)-\eta i\pi ae^{-a},
 \quad a>0,\quad J_\eta(0)=1.\nonumber
\end{align}
Here $\eta=+1$ denotes the earlier $12$ convention; $-1$ is its
conjugate. The signs are boundary data, not a real-part projection.

Put $A_r(\sigma)=\E[\operatorname{sech}^rX]$. The correctly normalized
four-time result is
\begin{equation}
 F_{{\rm disc},\sigma}=\frac{A_1^2}{16\alpha^2},\qquad
 {\cal R}_{\eta,\sigma}(t)
  =\frac{\E[C^{-1}R_\eta(t+Z,C)]}{A_1^2}.
 \label{smear}
\end{equation}
Both the random $C$ and its weight are necessary.

\newpage
\section{Convergence and a surviving contour difference}
The real-axis boundary kernels have a uniform bound $|J_\eta(c)|<6$.
For $c>0$ this follows from $0<J(c)\le1$.
For $c<0$ use $a=-1/c$.
If $a\le1$, the series
$\operatorname{Ei}(a)=\gamma_E+\log a+\sum_{k\ge1}a^k/(k\,k!)$
bounds the real part by $\gamma_E+1/e+e-1$.
If $a\ge1$, write
\[
 ae^{-a}\operatorname{Ei}(a)
 =ae^{-a}\operatorname{Ei}(1)
   +a\int_0^{a-1}\frac{e^{-v}}{a-v}\,dv .
\]
Splitting at $a/2$ bounds this by
$\operatorname{Ei}(1)/e+2+2/e$; the second interval is absent when
$a<2$. The imaginary magnitude is at most $\pi/e$.
Thus $|R_\eta|<13$ and Gaussian integration is absolutely convergent.

For standardized truncations $|X|,|Y|\le L\sigma/\sqrt2$,
$|Z|\le L\sigma$, the deleted contribution to the normalized numerator
is at most
\[
 \frac{39}{A_1^2}\operatorname{erfc}(L/\sqrt2).
 \label{tail}
\]
This is a continuum Gaussian-tail bound, not a Gauss--Hermite
quadrature error bound. Also
$A_1\ge e^{\sigma^2/4}\operatorname{erfc}(\sigma/2)>0$,
using $\operatorname{sech}x\ge e^{-|x|}$.

Suppose $\|h\|_\infty\le4$, a sufficient heavy-observer condition.
On the negative-$h$ interval let $a=4C/|h|\ge1$. Its imaginary
contribution to $R_+$ is
\[
 f_a(T)=\pi\left[ae^T e^{-ae^T}-ae^{-T}e^{-ae^{-T}}\right].
\]
It is odd, and strictly negative for $T>0$, since the ratio of
the first positive term to the second is
$\exp[2T-2a\sinh T]<1$.
Convolution with the centered Gaussian preserves this strict sign:
for $t>0$, pair $T>0$ and $-T$ and use
$p_\sigma(T-t)>p_\sigma(T+t)$.
The negative-$h$ angle set has positive measure for every
$\gamma\in[0,1)$. Averaging with $C^{-1}>0$ therefore proves
\begin{equation}
 \Im{\cal R}_{+,\sigma}(t)<0\quad(t>0),\qquad
 {\cal R}_{-,\sigma}(t)=\overline{{\cal R}_{+,\sigma}(t)},\qquad
 {\cal R}_{+,\sigma}(-t)=\overline{{\cal R}_{+,\sigma}(t)} .
 \label{sign}
\end{equation}
The unsmeared $\sigma=0$ limit has the same sign.

For Hermitian $W_\sigma,V_\sigma$, a genuine trace expression
$\operatorname{Tr}(yW_\sigma(t)yV_\sigma yW_\sigma(t)yV_\sigma)$
with $y=e^{-\beta H/4}$ is real whenever the products are trace class:
take its adjoint and use cyclicity. Equation \eqref{sign} shows that
real smearing alone does not turn either specified eikonal contour
into that expression. It does not exclude missing gravitational
contributions or alternative physical state prescriptions.

\newpage
\section{Fixed-order response and noncommuting limits}
Using the usual fixed-order asymptotic expansion
$J(c)\sim\sum_{n\ge0}n!(-c)^n$ \cite{dlmf}, \eqref{smear} gives
\begin{equation}
 {\cal R}_{\eta,\sigma}(t)\sim
 1+\sum_{n\ge1}\frac{n!(-1)^n\langle h^n\rangle}{4^n}
 Q_n(\sigma)(e^{nt}+e^{-nt}),\qquad
 Q_n=e^{n^2\sigma^2/2}\left(\frac{A_{n+1}}{A_1}\right)^2 .
 \label{pert}
\end{equation}
All coefficients are real; the two boundary values agree to every
fixed perturbative order. This is not a convergent series.

For fixed $\gamma,\sigma,t$ and $m\to\infty$, the fixed-order expansion
is justified by splitting $|Z|\le\tfrac12\log m$ from its complement.
On the inner region $|c_\pm|=O(m^{-1/2})$, the fixed-order $E_1/Ei$
remainder is bounded by a constant times $|c_\pm|^{N+1}$.
Its Gaussian expectation is $O(m^{-N-1})$.
The outer exact integrand is uniformly bounded; its Gaussian
probability, and the tails of any fixed polynomial approximation,
decay faster than every inverse power of $m$.
This argument is not uniform in $\sigma$, $t$ or expansion order.

A small-width expansion yields
\[
 A_r=1-\frac r4\sigma^2+\frac{r(3r+2)}{32}\sigma^4+O(\sigma^6),
 \quad
 \log Q_n=\frac{n(n-1)}2\sigma^2+\frac{n(n+3)}8\sigma^4+O(\sigma^6).
\]
In particular $Q_1=1+\sigma^4/2+O(\sigma^6)$: there is no
order-$\sigma^2$ change of the first gravitational coefficient.
For $Q_2$, the order-$\sigma^2$ term survives.
The first two terms are
\[
 1-\frac{\langle h\rangle}{4}Q_1(e^t+e^{-t})
   +\frac{\langle h^2\rangle}{8}Q_2(e^{2t}+e^{-2t}),
\]
where the existing shock moments are
\[
 \langle h\rangle=-\frac{4\pi G}{\alpha},\qquad
 \langle h^2\rangle=\frac{1+\sin\alpha/\alpha}
                     {8m^2\operatorname{sinc}^2(\pi\gamma)} .
\]
For pure recoil $\gamma=0$, the odd moments vanish. Neither
Gaussian smearing nor omission of this recoil term is a change of
the actual gravitational coupling.

If instead $t=\log m+u$ with fixed $u$, bounded convergence gives
\[
 {\cal R}_{\eta,\sigma}^{\rm boost}(u)
 =\frac1{A_1^2}\E\left[C^{-1}
  \left\langle J_\eta\left(\frac{H_\gamma e^{u+Z}}{4C}\right)
  \right\rangle\right].
\]
The opposite sector tends to one and cancels the subtraction.
The fixed-time limit is one; the boosted limit remains nontrivial.

These are statements about the smeared eikonal formula. Gaussian
tails sample arbitrarily large relative times. Without a uniform
bound on the omitted gravitational terms, the continuum-tail estimate
does not certify the difference from the true gravity correlator.

\newpage
\section{Numerical checks}
The main computation uses $m=1000$, $\gamma=0,0.25$,
$\sigma=0.4,0.8$, and $t=\log1000+u$, $u=0,2$.
We retain complex values of both prescriptions through conjugation,
without averaging them. Gauss--Hermite orders 20 and 32 use
24 and 40 Gauss--Legendre nodes per shock-sign interval.
A fresh order-48/64 run checks all eight points.
The maximum 20/32 change is $1.89\,10^{-7}$; the maximum
32/48 change is $1.38\,10^{-9}$. These are empirical changes,
not interval-certified errors.
\begin{center}
\begin{tabular}{rrrrr}
\toprule
$\gamma$ & $\sigma$ & $u$ & $\Re{\cal R}_{+,\sigma}$ &
 $\Im{\cal R}_{+,\sigma}$\\
\midrule
0 & .4 & 0 & 1.02238240 & -.00609498\\
0 & .8 & 0 & 1.02340116 & -.02263423\\
.25 & .4 & 0 & 1.06974804 & -.01230282\\
.25 & .8 & 0 & 1.06817688 & -.03757875\\
.25 & .8 & 2 & .94334065 & -.41973792\\
\bottomrule
\end{tabular}
\end{center}
At $(\gamma,\sigma,u)=(.25,.8,0)$, the incorrect center-only
convolution gives approximately $1.08131888-.05542995i$.
This is a measurable change in both parts, not just a choice of
notation for the smearing width.

An independent implementation retains the original four $d_i$,
constructs $A,B,P$ from the complex insertion coordinates, and uses
16/24-order four-dimensional Gaussian quadrature with 48/72 angular
nodes. At $u=0,\gamma=.25$, its fine values differ from the reduced
order-32 result by $2.03\,10^{-14}$ for $\sigma=.4$ and
$2.16\,10^{-9}$ for $\sigma=.8$.
This tests the change of variables and disconnected weight.

Additional controls compare the original one-dimensional (G.2)
momentum integral with the reduced expression at three positive-$h$
configurations using 70 decimal digits; check 13 positive, negative
and small-$c$ values against mpmath; and independently integrate
the $A_r$ moments. The 52 diagnostic criteria all pass.
A further 11 checks cover the eight refinements, exact Gaussian
covariance, the small-width polynomial, and retained diagnostics.
\begin{center}
\begin{tabular}{rrr}
\toprule
$\sigma$ & $Q_1$ & $Q_2$\\
\midrule
.1 & 1.00004927 & 1.01017421\\
.2 & 1.00075574 & 1.04275682\\
.4 & 1.01049126 & 1.20343051\\
.8 & 1.11849529 & 2.45158736\\
\bottomrule
\end{tabular}
\end{center}
SciPy double precision uses explicit lateral imaginary parts and a
24-term small-$c$ expansion for $|c|<1/128$ to avoid overflow.
No tolerance or saved reference was changed after the checks.
Parameters, dependency versions, raw runs, source hashes and
compilation evidence are in \path{results/stage4/gaussian-eikonal/}.

\newpage
\section{What this resolves, and what remains}
The independent insertions now have an explicit finite-width
eikonal comparison formula, a correctly smeared disconnected
baseline, and a sign test for the remaining contour difference.
This is a self-checked consequence of established formulas, with
no claim of priority for the smearing method or the underlying
reality obstruction.

The conformal free antipodal two-point function is
$[4\alpha\cosh(\tau/2)]^{-1}$. Smearing a pair gives
$A_1/(4\alpha)$, explaining the disconnected baseline directly.
The same real time-smearing prescription may be applied to the
Gaussian-regulated free field before the antipodal continuation.
This matches the \emph{protocol}; it does not identify the fixed-mass
eikonal state with the preceding principal-series Fock vacuum,
or extend that theorem to complementary-series matter.

A substantive next comparison must keep the actual real-principal-series
probe masses, the quantum observer preparation, and all vacuum
contractions. For general probes the Bessel functions in CSTY (2.32)
replace the conformal exponential. Their two-point factors depend on
the independent pair differences. A direct smearing of a normalized
symmetric-time curve will again discard this information.

For a microscopic comparison, one also needs the physical finite-parameter
insertion normalization and the appropriate inner product/Gram inverse,
with its analytic-continuation contour. A positive Euclidean graph
amplitude alone does not fix either. The sign result here isolates
a question for such a construction; it is not resolved by taking the
real part of one of the two candidate contours.

New adjacent sources located during the targeted literature check are
the observer horizon-edge-mode EFT \cite{edge} and the graviton
eikonal analysis \cite{mnx}. Only their abstracts and introductory
scope were read this round. They should next be examined for their
treatment of lifted observer modes, zero modes and state preparation;
they are not inputs to the derivation above.

\begin{thebibliography}{9}
\bibitem{csty} Y. Chen, D. Stanford, H. Tang and Z. Yang,
\href{https://arxiv.org/html/2607.14042v1}
{Negative shocks versus static patch holography},
Eqs.\ (2.25)--(2.26), (2.32)--(2.42), Appendix F and (G.2).
\bibitem{clp} J. Chandra, D. Liska and B. Post,
\href{https://arxiv.org/html/2610.30247v1}
{3D de Sitter from complex Virasoro TQFT: Observers and knots},
Eqs.\ (3.35)--(3.38) and Section 6.
\bibitem{dlmf} NIST Digital Library of Mathematical Functions,
\href{https://dlmf.nist.gov/6.12}{Section 6.12(i)},
exponential-integral asymptotics and fixed-order remainder bounds.
\bibitem{edge} J. Bonifacio, K. Hinterbichler, Y. T. A. Law and V. Lochab,
\href{https://arxiv.org/html/2610.28719v1}
{De Sitter Observers and an Effective Field Theory of Horizon Edge Modes}.
\bibitem{mnx} A. Milekhin, V. Narovlansky and J. Xu,
\href{https://arxiv.org/html/2607.13137v1}
{Out-of-time-ordered Correlators in de Sitter Revisited}.
\end{thebibliography}
\end{document}
